% Research draft: AI-assisted derivation and internal mathematical review.
% Not a human-specialist validation or historical-priority claim.
\begin{theorem}[A dimension-dependent batch guarantee]
\label{thm:helly-precision}
Let $F$ be nonempty, compact and convex, and let $T(x)=c^\top x+d$.
Fix a first-stage exact linear observation $L_0x$, whose feasible-answer
set $I=L_0F$ has affine dimension $r$. After this answer, a deterministic
policy chooses one of finitely many second-stage linear measurement
batches $L_1,\ldots,L_m$. Every combined batch $(L_0,L_j)$ must obey the
same declared hard budget, and the eligible second-stage batches must not
depend on the first answer. Suppose the policy's terminal sharp interval
has width at most $\tau$ for every $x\in F$.
Then some randomized nonadaptive policy supported on at most $r+1$
of these combined batches has worst-case precision success at least
$1/(r+1)$.
In particular, a sure-precision policy using two scalar queries yields
a randomized-batch success guarantee of at least $1/2$, even when its
conditional worst widths are not affine functions of the first answer.
\end{theorem}
\begin{proof}
For $t\in I$, define the conditional worst terminal width
\[
 g_j(t)=\max\{c^\top(u-v):u,v\in F,\ L_0u=L_0v=t,
                                      \ L_ju=L_jv\}.
\]
The feasible set is nonempty and compact, so the maximum is attained.
The expression is exactly the largest terminal sharp width among records
with first answer $t$: an endpoint pair with equal first and second
answers belongs to a common terminal fiber, and the endpoints of every
such fiber are feasible in this maximization.

Each $g_j$ is concave on $I$. Indeed, convex combinations of optimizing
endpoint pairs at $t_1,t_2$ are feasible at the corresponding convex
combination of answers, and attain the same convex combination of their
objective values. The function is also upper semicontinuous: for a
convergent sequence of answers, compactness gives convergent subsequences
of optimizing endpoints, whose limits are feasible at the limiting
answer. Consequently, $g_j$ is finite and measurable.

For a fixed first answer, the policy chooses a single second-stage batch;
its sure guarantee therefore implies $\min_jg_j(t)\leq\tau$ throughout
$I$. The strict failure sets
\[
 H_j=\{t\in I:g_j(t)>\tau\}
\]
are convex by concavity, and $\bigcap_{j=1}^m H_j=\varnothing$.
Finite Helly's theorem, applied in the affine hull of $I$, gives a
subfamily $J$ of at most $r+1$ sets with empty intersection. This finite
form of Helly's theorem does not require the sets to be closed; thus
using strict failure sets preserves success at equality with $\tau$.

Choose uniformly among the combined batches indexed by $J$. For every
fixed record $x$, its answer $t=L_0x$ lies outside at least one $H_j$,
so that batch's actual terminal width is at most $g_j(t)\leq\tau$.
Its precision success probability is therefore at least $1/|J|$, and
hence at least $1/(r+1)$, for every record.
\end{proof}

\paragraph{Stopping and admissibility.}
Stopping after the first stage is included by allowing the empty
second-stage batch. Stopping without any query is a fixed batch and
already supplies success one when it guarantees the precision.
For ordinary two-query policies, the deterministic first query is fixed
before any answer, and $r\leq1$. Unequal fixed query costs are allowed
provided that every batch used in the support is admissible under the
common hard budget. The result does not allow an answer-dependent query
dictionary, an answer-dependent budget, or errors chosen after seeing
the algorithm's random seed.

\paragraph{Sharpness.}
The manuscript's existing $K$-capacity construction has first-stage
answer dimension $r=K-1$, using its $K-1$ explicit capacities. Its second
stage reads all but one of the $K$ bounded quantities. The adaptive policy
guarantees width $1/K$, whereas even unrestricted admissible randomized
batches have worst-case precision success exactly $1/K=1/(r+1)$.
Thus the dimension-dependent constant is sharp, using the already proved
family. For $r=1$, the three-coordinate example is sharp.

\paragraph{Extraction on polytopes.}
For a scalar first answer and a guaranteed threshold, at most two
second-stage actions suffice. A pair $(j,k)$ can be checked with one LP:
give the two actions separate endpoint pairs sharing a first answer $t$,
maximize $\eta$, and require both endpoint contrasts to be at least
$\eta$. Its value is $\max_{t\in I}\min\{g_j(t),g_k(t)\}$.
The pair covers every first answer precisely when this value is at most
$\tau$. Testing pairs, together with singleton actions, constructs the
nonadaptive support without gridding possible answers. This uses at most
$\binom{m}{2}+m$ LPs after the first stage has been chosen.

The same argument implies
\[
 \max_{t\in I}\min_{1\leq j\leq m}g_j(t)
 =\min_{\substack{J\subseteq\{1,\ldots,m\}\\1\leq|J|\leq2}}
     \max_{t\in I}\min_{j\in J}g_j(t)
\]
for scalar first answers. Compactness and upper semicontinuity give an
attained left-hand maximum, so the Helly argument applies at its exact
value. This is a sparse sufficient policy, not a claim that the same two
actions minimize conditional width at every answer.
